\section{Human-in-the-Loop Implementation}
\label{sec:hitl_impl}

The framework places humans at critical decision points through LangGraph's \texttt{interrupt()} primitive, which pauses execution and returns control to the API server until the user responds.

\subsection{Interrupt Mechanism}

LangGraph checkpoints the full graph state after every node execution. When a node calls \texttt{interrupt()}, it suspends and the runtime surfaces the payload to the calling client. The code example below, from the firmware-to-AI transition gate, shows the pattern:

\begin{lstlisting}[language=Python, caption={Interrupt-based approval gate}]
# From workflow1_firmware.py
def decide_continue_to_ai(state, config):
    # Suspend execution and request user input
    user_response = interrupt({
        "type": "decision_request",
        "message": "Firmware generated. Continue with AI model selection?",
        "options": ["yes", "no", "skip"]
    })
    
    # User responses are classified via LLM
    llm_parser = get_llm(config, structured_schema=DecisionExtraction, temperature=0)
    result = llm_parser.invoke([
        SystemMessage(content=decision_extraction_instructions),
        HumanMessage(content=user_response)
    ])
    
    if result.continue_to_ai:
        state.route = "ai_flow"
    else:
        state.route = "END"
    
    return state
\end{lstlisting}

When the client sends the user's reply, the graph resumes from the exact checkpoint, with all state variables intact.

\subsection{Approval Gates vs. Steering Inputs}

The framework distinguishes two HITL patterns.

\subsubsection{Human-as-Gate}

Before any high-risk operation, the workflow suspends and requests an explicit approve or reject command. In practice this covers three points: confirming the transition from firmware generation to AI model selection, accepting or rejecting a candidate model during search, and approving the fine-tuned model before deployment after NNI completion.

\subsubsection{Human-as-Instructor}

Natural language steering lets the user redirect the workflow mid-execution without predefined menu options:

\begin{quote}
    User: \textit{``Actually, I want a smaller model''}
\end{quote}

The response is parsed with structured extraction and the state is updated accordingly:

\begin{lstlisting}[language=Python]
result = llm_classifier.invoke(user_input)
# result.intent = "change_model_size"
# result.target_size = "small"

state.model_size_preference = "small"
state.route = "search_model"  # Return to model discovery
\end{lstlisting}

\subsection{State Persistence and Resume Logic}

LangGraph stores checkpoints in \textbf{Redis} using a hierarchical key structure. Each entry is serialized as a MessagePack or JSON payload:

\begin{verbatim}
checkpoint:{thread_id}:{checkpoint_id} -> {
    "v": 1,
    "ts": "2023-10-27T10:00:00Z",
    "channel_values": {
        "messages": [...],
        "firmware_project_dir": "/path/to/project",
        "route": "ai_flow"
    },
    "parent_config": { ... }
}
\end{verbatim}

Resume begins when the client posts the user's reply to \texttt{POST /threads/\{thread\_id\}/runs}. LangGraph retrieves the latest checkpoint for the thread, deserializes \texttt{MasterState}, injects the new input, and resumes from the next node. This mechanism supports multi-hour workflows (overnight NNI experiments, for instance) without any explicit session management code.

\subsection{Actionable Reporting and Synthetic Summaries}

Raw toolchain outputs from CubeMX or STEdgeAI can run to hundreds of lines. Rather than exposing this verbatim to the user, the orchestrator applies heuristic log parsing after each tool execution and extracts a single-line summary: for example, \textit{``Project generated successfully (STM32H7, 4 peripherals enabled)''} or \textit{``WARNING: Memory Overflow: Flash required 3.2MB/2MB''}. These summaries are stored in \texttt{MasterState} and used to build context-aware prompts at interrupt points, reducing the information the user needs to review before making a decision.
